% \documentclass[11pt, a4paper]{article}

% % --- Encodage et Langue ---
% \usepackage[utf8]{inputenc}
% \usepackage[T1]{fontenc}
% \usepackage[french]{babel}

% % --- Mise en page et Graphismes ---
% \usepackage[margin=2cm]{geometry}
% \usepackage{pdflscape} % Pour le tableau en mode paysage
% \usepackage[table]{xcolor}
% \usepackage{float}
% \usepackage{booktabs}
% \usepackage{array} % Pour personnaliser les colonnes tabular
% \usepackage{hyperref}

% % --- Mathématiques ---
% \usepackage{amsmath, amssymb, amsthm}

% % --- Couleurs personnalisées ---
% \definecolor{headerblue}{RGB}{31, 73, 125}
% \definecolor{rowgray}{RGB}{242, 242, 242}
% \definecolor{rowwhite}{RGB}{255, 255, 255}

% \title{\textbf{Note de Synthèse : Modélisation du Football}}
% \author{Stage M1 - Mathématiques Appliquées}
% \date{\today}

% \begin{document}

% \maketitle
% \tableofcontents
% \newpage

% \section{Dixon \& Coles (1997)}
% L'article de Dixon et Coles \cite{dixon} est une référence séminale qui adapte le modèle de Maher (1982). Ils utilisent une approche fréquentiste basée sur la \textbf{loi de Poisson} pour estimer les capacités d'attaque et de défense de chaque équipe.

% \subsection{Données et Inflexion}
% L'étude exploite un historique de 6 629 matchs du championnat anglais (1992-1995). Contrairement aux modèles antérieurs, ils introduisent deux innovations majeures :
% \begin{itemize}
%     \item \textbf{La dépendance des scores ($\rho$)} : Correction pour les scores faibles (0-0, 1-1) qui ne suivent pas une loi de Poisson pure.
%     \item \textbf{Le dynamisme temporel} : Les performances passées sont pondérées pour accorder plus d'importance aux résultats récents.
% \end{itemize}

% \section{Egidi \& Torelli (2021)}
% Cet article \cite{egidi} déplace le curseur vers la méthodologie statistique. Ils comparent deux philosophies : prédire le nombre de buts (\textit{Goal-based}) ou prédire l'issue du match (\textit{Result-based}).

% L'approche est ici \textbf{Bayésienne}, utilisant des simulations MCMC pour capturer l'incertitude des paramètres de compétence des équipes, testée sur les données de la Coupe du Monde 2018.

% \section{Carloni et al. (2021)}
% Le travail de Luca Carloni et son équipe \cite{carloni} représente le versant \textbf{Intelligence Artificielle}. L'objectif n'est plus d'expliquer mathématiquement le jeu, mais de maximiser le \textbf{ROI (Retour sur Investissement)}. Ils utilisent des réseaux de neurones (ANN) entraînés sur une base de données massive issue du web scraping pour battre les cotes des bookmakers.

% \newpage

% % ================================================================
% % TABLEAU COMPARATIF EN PAYSAGE
% % ================================================================
% \begin{landscape}
% \section{Synthèse comparative des approches}

% \renewcommand{\arraystretch}{1.8}
% \begin{table}[H]
% \centering
% \caption{Grille d'analyse des trois piliers de la littérature étudiée}
% \small
% \begin{tabular}{|>{\bfseries}p{3.5cm}|p{6cm}|p{6cm}|p{6cm}|}
% \hline
% \rowcolor{headerblue} 
% \textcolor{white}{Critères} & 
% \textcolor{white}{Dixon \& Coles (1997)} & 
% \textcolor{white}{Egidi \& Torelli (2021)} & 
% \textcolor{white}{Carloni et al. (2021)} \\ \hline

% \rowcolor{rowgray}
% Approche & Statistique fréquentiste & Statistique Bayésienne & Machine Learning (IA) \\ \hline

% \rowcolor{rowwhite}
% Modèle principal & Poisson bivarié ajusté & Inférence Bayésienne (MCMC) & Réseaux de neurones (ANN) \\ \hline

% \rowcolor{rowgray}
% Cible de prédiction & Nombre de buts (HG, AG) & Comparaison Buts vs Issue & Issue du match (1N2) \\ \hline

% \rowcolor{rowwhite}
% Innovation clé & Dépendance $\rho$ et facteur temporel $\alpha$ & Diagnostic de calibration (Brier Score) & Approche "Data-driven" / ROI \\ \hline

% \rowcolor{rowgray}
% Points forts & Très interprétable, robuste & Gestion fine de l'incertitude & Capture les relations non-linéaires \\ \hline

% \rowcolor{rowwhite}
% Points faibles & Ignore le contexte (blessures, etc.) & Calculs lourds (simulations) & Effet "Boîte noire" \\ \hline
% \end{tabular}
% \end{table}
% \end{landscape}

% \newpage
% \nocite{*}
% \bibliographystyle{plain}
% \bibliography{references}

% \end{document}



\documentclass{article}
\usepackage{lscape} % Pour mettre le tableau en paysage
\usepackage[utf8]{inputenc}
\usepackage[T1]{fontenc}
\usepackage[french]{babel}
\usepackage{array,multirow,makecell}
\setcellgapes{1pt}
\makegapedcells
\newcolumntype{R}[1]{>{\raggedleft\arraybackslash }b{#1}}
\newcolumntype{L}[1]{>{\raggedright\arraybackslash }b{#1}}
\newcolumntype{C}[1]{>{\centering\arraybackslash }b{#1}}

\begin{document}

% \begin{tabular}{|R{2cm}||C{2cm}|C{2cm}|C{2cm}|}
% \hline  & Dixon et Coles &  Egidi et Torelli &  Carloni et al.  \\
% \hline type de données & xx & xx & xx  \\
% \hline hypothèses mathématiques & xx & xx & xx \\
% \hline métriques d'évaluation & xx & xx & xx  \\
% \hline
% \end{tabular}
\begin{landscape}

\begin{tabular}{|L{2cm}||C{5cm}|C{5cm}|C{5cm}|} 
    \hline  & \textbf{Dixon et Coles} &  \textbf{Egidi et Torelli} &  \textbf{Carloni et al.}  \\
    \hline \textbf{Type de données} & Résultats historiques de 6 629 matchs (scores) de la ligue et coupes anglaises (1992-1995) 
    & matchs de la Coupe du Monde FIFA 2018, incluant les classements FIFA et les cotes de bookmakers
    & Données scrapées de 49 319 matchs (12 pays, 2008-2020) incluant les cotes de paris (1X2, Over-Under) \\
    \hline \textbf{Hypothèses mathématiques} & Lois de Poisson pour les buts avec paramètres attaque/défense, 
    correction de dépendance ($\tau$) pour les petits scores et pondération temporelle ($\Phi$)
    & Cadre bayésien comparant des modèles multinomiaux (résultats) et des modèles de Poisson (buts) basés sur le différentiel de classement FIFA
    & Utilisation d'algorithmes de Machine Learning (SVM, Random Forest) avec une priorité donnée au Réseau de Neurones Artificiels (ANN) \\
    \hline \textbf{Métriques d'évaluation} & Profitabilité et retour positif d'une stratégie de paris, ainsi 
    que la vraisemblance prédictive 
    & Pseudo-$R^2$ , Brier Score (erreur quadratique moyenne) et Leave-One-Out Cross Validation (LOOIC)
    & Précision de la prédiction (Accuracy) et Retour sur Investissement (ROI/RSI) \\ 
    \hline 
\end{tabular}

\end{landscape}


\end{document}